% ============================================================
% SECTION VI: DISCUSSION + CONCLUSION
% ============================================================
\section{Discussion}

\textbf{Two-token attention degeneracy.}
The observation that self-attention over a length-2 sequence boils down to scalar interpolation is not specific to deepfake forensics. Any dual-branch architecture that fuses pooled embeddings via self-attention should watch for this failure mode. With $h = 8$ heads, attention yields 8 scalar weights; our gating yields 512 per-dimension weights. We suspect this insight applies to other dual-stream tasks as well - for example, audio-visual fusion - though validating that is left for future work.

\textbf{Asymmetric compute placement.}
Running SRM on CPU (inside DataLoader workers) and FFT on GPU is a practical design choice that tends to get overlooked in architecture papers. It avoids two bottlenecks: doing both in the forward pass starves the GPU of compute; pre-computing both to disk costs around 240\,GB for the full FF++ dataset.

\textbf{Identity-safe evaluation.}
The 2-5\% accuracy drop from identity-safe splits is an honesty measure, not a flaw. We would encourage the community to adopt source-ID-based partitioning as the default for attribution benchmarks.

\subsection{Limitations}
The pipeline tracks a single face per video. Training covers four FF++ methods; newer generators (diffusion-based) are not represented. FF++ is a controlled-lab dataset, and the cross-dataset generalization gap remains substantial. Heavy compression (c40) may degrade the SRM/FFT signals that the frequency stream relies on.

\section{Conclusion}

We presented DSAN, a dual-stream attribution network that goes beyond asking \textit{whether} facial media has been manipulated to identify \textit{how}. The gated bilinear fusion mechanism solves a concrete problem - the degeneracy of two-token self-attention - and lets the frequency stream meaningfully contribute to attribution decisions. Multi-task training with cross-entropy, supervised contrastive learning, and blending-mask prediction gives the model complementary learning signals, while SBI augmentation helps with cross-dataset robustness. The dual Grad-CAM++ explainability surfaces method-specific saliency in both spatial and frequency domains. Taken together, the modular forensic pipeline shows that deepfake investigation can be turned into a reproducible, auditable process.

\textbf{Future work} includes extending the system to diffusion-based generators, running full-scale cross-dataset evaluations, testing robustness under c40 compression, and exploring cross-attention fusion as an alternative to gating.

\section*{Acknowledgment}
The author thanks Dr.\ Kaptan Singh for supervision and guidance. The FaceForensics++ dataset was used under academic license.
